\documentclass[12pt]{article}
\usepackage{amsmath,amssymb,amsfonts,amsthm}
\usepackage[margin=1in]{geometry}

\begin{document}

\begin{center}
{\Large\bfseries Chapter 9, Problem 9}\\[6pt]
Byron \& Fuller, \textit{Mathematics of Classical and Quantum Physics}
\end{center}

\bigskip

\textbf{Problem.} Consider the hydrogen atom with Hamiltonian $H_0$ and eigenstates $\psi_{n\ell m}$. The $2s$--$2p$ states are nearly degenerate (split by the Lamb shift $L \approx 4 \times 10^{-6}$ eV). An electric field $H_1 = eEz$ is applied.

(a) Explain why we can restrict attention to the $2s$--$2p$ subspace.

(b) Show that the Brillouin--Wigner perturbation series (Eq.~9.34) has nonvanishing first two terms for $\psi_{2s}$ in terms of $\psi_{2p_0}$, while the Rayleigh--Schr\"odinger series (Eq.~9.40) would not.

(c) Find $\psi_{2s}$ and $\lambda_{2s}$ in terms of $L$ and $M = \langle\psi_{2s}, H_1\psi_{2p_0}\rangle$.

(d) Take the limit $L \to 0$.

\bigskip
\textbf{Solution.}

\subsection*{(a) Restriction to the $2s$--$2p$ subspace}

The perturbation $H_1 = eEz$ shifts energies by amounts of order $M \sim eEa_0 \sim L$ (where $a_0$ is the Bohr radius). The Lamb shift $L \approx 4 \times 10^{-6}$ eV separates the $2s$ and $2p$ levels.

The nearest states outside the $2s$--$2p$ manifold are the $3s$, $3p$, $3d$ states, which lie about $\Delta E \approx 2$ eV away. Since $L \ll \Delta E$, the mixing with these distant states is of order $M/\Delta E \ll 1$ and can be neglected. The perturbation theory of Section~9.2 (for nearly degenerate states) is specifically designed for this situation: states that are nearly degenerate must be treated non-perturbatively among themselves, while far-away states contribute negligible corrections.

Furthermore, by the selection rules for the matrix elements of $z = r\cos\theta$, $H_1$ only connects states whose $\ell$ values differ by 1, so within the $2s$--$2p$ manifold, only $\langle 2s|H_1|2p_0\rangle = M$ is nonzero (the $2p_{\pm}$ states have $\Delta m = \pm 1$ and don't couple to $2s$ via $z$).

\subsection*{(b) Perturbation series terms}

Let $\lambda_{2s}^{(0)}$ and $\lambda_{2p}^{(0)} = \lambda_{2s}^{(0)} - L$ be the unperturbed energies (with $2s$ above $2p$ by $L$).

\textbf{Brillouin--Wigner (Eq.~9.34):} The perturbed eigenvalue $\lambda_{2s}$ appears in the denominators. The expansion involves:
\[
\psi_{2s} = \psi_{2s}^{(0)} + \frac{\langle\psi_{2p_0}, H_1\psi_{2s}^{(0)}\rangle}{\lambda_{2s} - \lambda_{2p}^{(0)}}\,\psi_{2p_0} + \cdots
\]
The first correction term involves $M^*/(\lambda_{2s} - \lambda_{2p}^{(0)})$, which is nonzero since $M \neq 0$ and the denominator involves the (unknown) perturbed energy, not zero. The second-order term also contributes.

\textbf{Rayleigh--Schr\"odinger (Eq.~9.40):} Here the denominator is $\lambda_{2s}^{(0)} - \lambda_{2p}^{(0)} = L$. Since $L$ is of the same order as the perturbation matrix element $M$, the RS series does not converge properly---the ``small'' expansion parameter $M/L$ is of order unity, so the first few terms of the RS series are all comparable in size and the series is not useful. The BW series, by using the exact energy in the denominator, partially resums the RS series and gives a more useful result.

\subsection*{(c) Finding $\psi_{2s}$ and $\lambda_{2s}$}

Within the $2s$--$2p_0$ subspace (ignoring $2p_{\pm}$ since they don't couple), the effective Hamiltonian matrix is:
\[
H_{\text{eff}} = \begin{pmatrix} \lambda_{2s}^{(0)} & M \\ M^* & \lambda_{2p}^{(0)} \end{pmatrix} = \begin{pmatrix} E_0 + L & M \\ M^* & E_0 \end{pmatrix},
\]
where $E_0 = \lambda_{2p}^{(0)}$ and we take $M$ to be real for simplicity.

The eigenvalue equation $\det(H_{\text{eff}} - \lambda I) = 0$ gives:
\[
(E_0 + L - \lambda)(E_0 - \lambda) - M^2 = 0.
\]
\[
\lambda^2 - (2E_0 + L)\lambda + E_0(E_0 + L) - M^2 = 0.
\]
\[
\lambda = E_0 + \frac{L}{2} \pm \frac{1}{2}\sqrt{L^2 + 4M^2}.
\]

The state that reduces to $\psi_{2s}$ when $E \to 0$ (i.e., $M \to 0$) corresponds to the $+$ sign:
\[
\boxed{\lambda_{2s} = E_0 + \frac{L}{2} + \frac{1}{2}\sqrt{L^2 + 4M^2}}.
\]

The corresponding eigenvector (unnormalized) is:
\[
\psi_{2s} \propto \psi_{2s}^{(0)} + \frac{M}{\lambda_{2s} - E_0}\,\psi_{2p_0} = \psi_{2s}^{(0)} + \frac{2M}{L + \sqrt{L^2 + 4M^2}}\,\psi_{2p_0}.
\]

Using the Brillouin--Wigner form (Eqs.~9.33--9.34):
\[
\boxed{\psi_{2s} = \mathcal{N}\left[\psi_{2s}^{(0)} + \frac{2M}{L + \sqrt{L^2 + 4M^2}}\,\psi_{2p_0}\right]},
\]
where $\mathcal{N}$ is the normalization constant.

\subsection*{(d) Limit $L \to 0$}

When $L \to 0$ (exact degeneracy):
\[
\lambda_{2s} \to E_0 + |M|, \qquad \sqrt{L^2 + 4M^2} \to 2|M|.
\]

The coefficient of $\psi_{2p_0}$ becomes $\frac{2M}{0 + 2|M|} = \text{sgn}(M) = \pm 1$ (assuming $M > 0$, this is $+1$).

So:
\[
\psi_{2s} \to \frac{1}{\sqrt{2}}(\psi_{2s}^{(0)} + \psi_{2p_0}), \qquad \lambda_{2s} \to E_0 + M.
\]

This is the standard result for degenerate perturbation theory: the electric field creates a \emph{linear Stark effect} with energy splitting $\pm M$, and the perturbed states are equal superpositions of $s$ and $p$ states. This is the well-known linear Stark effect of the hydrogen atom. \hfill $\blacksquare$

\end{document}
